% !TEX program = xelatex
% Lernplatt - by Can Siebert
% Kurs 71 (Dig 42) - Tag 16 - Lehrskript
% EU-Taxonomie & CSRD als Prozess-, Daten- und Kontrollarchitektur
%
% Self-contained xelatex source. Kein biblatex, keine externen Bilder.

\documentclass[11pt,a4paper,oneside]{article}

% --- Encoding & Sprache --------------------------------------------------
\usepackage{polyglossia}
\setdefaultlanguage{german}

% --- Geometrie -----------------------------------------------------------
\usepackage[a4paper,top=2cm,bottom=2cm,outer=2.5cm,inner=2.5cm,bindingoffset=1.5cm]{geometry}

% --- Fonts ---------------------------------------------------------------
\usepackage{fontspec}
\defaultfontfeatures{Ligatures=TeX}

\IfFontExistsTF{Charter}{%
  \setmainfont{Charter}[Numbers=OldStyle]%
}{%
  \IfFontExistsTF{Noto Serif}{%
    \setmainfont{Noto Serif}%
  }{%
    \setmainfont{Liberation Serif}%
  }%
}

\IfFontExistsTF{Source Sans 3}{%
  \setsansfont{Source Sans 3}%
}{%
  \IfFontExistsTF{Source Sans Pro}{%
    \setsansfont{Source Sans Pro}%
  }{%
    \setsansfont{Liberation Sans}%
  }%
}

\newcommand{\caveatfontfallback}{\itshape}
\IfFontExistsTF{Caveat}{%
  \newfontfamily\caveatfont{Caveat}[Scale=1.15]%
}{%
  \newcommand\caveatfont{\caveatfontfallback}%
}

\setmonofont{Liberation Mono}

% --- Mikro-Typografie ----------------------------------------------------
\usepackage{microtype}
\linespread{1.15}

% --- Farben & Boxen ------------------------------------------------------
\usepackage{xcolor}
\definecolor{lpInk}{RGB}{20,28,38}
\definecolor{kurs71}{HTML}{9D4A1A}
\definecolor{lpAccent}{HTML}{9D4A1A}
\definecolor{lpAccentLight}{RGB}{248,232,218}
\definecolor{lpAmber}{RGB}{222,138,30}
\definecolor{lpAmberLight}{RGB}{255,243,217}
\definecolor{lpAmberBorder}{RGB}{196,118,18}
\definecolor{lpGray}{RGB}{96,104,114}
\definecolor{lpGrayLight}{RGB}{238,240,243}

\usepackage[most]{tcolorbox}
\tcbuselibrary{breakable,skins}

\newtcolorbox{eselsbox}[1][Eselsbrücke]{%
  enhanced, breakable,
  colback=lpAmberLight,
  colframe=lpAmberBorder,
  boxrule=0.6pt,
  arc=2pt,
  borderline={0.4pt}{0pt}{lpAmberBorder, dashed},
  left=10pt,right=10pt,top=8pt,bottom=8pt,
  fonttitle=\sffamily\bfseries\color{lpAmberBorder},
  coltitle=lpAmberBorder,
  title={#1},
  attach boxed title to top left={xshift=8pt,yshift=-8pt},
  boxed title style={size=small, colback=lpAmberLight, colframe=lpAmberBorder, boxrule=0.4pt, arc=1pt},
  top=14pt
}

\newtcolorbox{merkbox}[1][Merksatz]{%
  enhanced, breakable,
  colback=lpAccentLight,
  colframe=lpAccent,
  boxrule=0.6pt,
  arc=2pt,
  left=10pt,right=10pt,top=8pt,bottom=8pt,
  fonttitle=\sffamily\bfseries\color{white},
  coltitle=white,
  title={#1},
  attach boxed title to top left={xshift=8pt,yshift=-8pt},
  boxed title style={size=small, colback=lpAccent, colframe=lpAccent, boxrule=0.4pt, arc=1pt},
  top=14pt
}

% --- Listen --------------------------------------------------------------
\usepackage{enumitem}
\setlist{itemsep=2pt, topsep=4pt, parsep=0pt}
\setlist[itemize,1]{label={\textbullet},leftmargin=1.6em}
\setlist[itemize,2]{label={\textendash},leftmargin=1.4em}
\setlist[enumerate,1]{label=\arabic*., leftmargin=1.8em}

% --- Tabellen ------------------------------------------------------------
\usepackage{tabularx}
\usepackage{array}
\usepackage{booktabs}
\usepackage{longtable}

% --- Kopf-/Fusszeilen ----------------------------------------------------
\usepackage{fancyhdr}
\usepackage{lastpage}
\pagestyle{fancy}
\fancyhf{}
\renewcommand{\headrulewidth}{0.3pt}
\renewcommand{\footrulewidth}{0pt}
\fancyhead[L]{\footnotesize\sffamily\color{lpGray}Lernplatt \textperiodcentered{} Kurs 71 \textperiodcentered{} Tag 16}
\fancyhead[R]{\footnotesize\sffamily\color{lpGray}CSRD \textperiodcentered{} EU-Taxonomie \textperiodcentered{} Controls}
\fancyfoot[L]{\footnotesize\sffamily\color{lpGray}\textcopyright{} Can Siebert}
\fancyfoot[R]{\footnotesize\sffamily\color{lpGray}\thepage{} / \pageref{LastPage}}

\fancypagestyle{plain}{%
  \fancyhf{}%
  \renewcommand{\headrulewidth}{0pt}%
  \fancyfoot[C]{\footnotesize\sffamily\color{lpGray}\thepage{} / \pageref{LastPage}}%
}

% --- Section-Styling -----------------------------------------------------
\usepackage[explicit]{titlesec}

\titleformat{\section}[block]
  {\sffamily\Large\bfseries\color{lpAccent}}
  {\thesection.}{0.6em}{#1}
  [{\vspace{2pt}\color{lpAccent}\titlerule[0.6pt]}]

\titleformat{\subsection}[block]
  {\sffamily\large\bfseries\color{lpInk}}
  {\thesubsection}{0.6em}{#1}

\titleformat{\subsubsection}[block]
  {\sffamily\normalsize\bfseries\color{lpInk}}
  {}{0pt}{#1}

\titlespacing*{\section}{0pt}{14pt}{8pt}
\titlespacing*{\subsection}{0pt}{10pt}{4pt}
\titlespacing*{\subsubsection}{0pt}{8pt}{2pt}

% --- Hyperlinks ----------------------------------------------------------
\usepackage[hidelinks,unicode]{hyperref}

% --- TikZ ----------------------------------------------------------------
\usepackage{tikz}
\usetikzlibrary{shapes.geometric,arrows.meta,positioning,calc,fit,backgrounds}
\definecolor{lpGreenLight}{RGB}{222,239,224}

% --- TOC -----------------------------------------------------------------
\renewcommand{\contentsname}{\sffamily\Large\bfseries\color{lpAccent}Inhalt}

% --- Helfer --------------------------------------------------------------
\newcommand{\akzent}[1]{{\caveatfont\color{lpAmber}#1}}
\newcommand{\fachbegriff}[1]{\textbf{#1}}
\newcommand{\quelle}[1]{{\footnotesize\color{lpGray}(#1)}}

% =========================================================================
\begin{document}

% --- COVER ---------------------------------------------------------------
\begin{titlepage}
  \pagestyle{empty}
  \vspace*{1.5cm}
  \noindent{\sffamily\footnotesize\color{lpGray}LERNPLATT \textperiodcentered{} BY CAN SIEBERT}\\[2pt]
  \noindent{\color{lpAccent}\rule{\linewidth}{1.2pt}}\\[4pt]
  \noindent{\sffamily\footnotesize\color{lpGray}MODUL 8 \textperiodcentered{} TAG 16 VON 16 \textperiodcentered{} KURS 71 (Dig 42) \textperiodcentered{} 19.\,Juni 2026}

  \vspace{3cm}
  {\sffamily\fontsize{30}{34}\selectfont\bfseries\color{lpInk}Regulatorik\\in Steuerung\\übersetzen\par}

  \vspace{0.7cm}
  {\sffamily\large\color{lpGray}EU-Taxonomie und CSRD als Prozess-, Daten- und Kontrollarchitektur --\\auditfähig und integriert in Transformation.\par}

  \vspace{1.5cm}
  {\caveatfont\LARGE\color{lpAmber}Lehrskript -- Tag 16 von 16}\par

  \vfill

  \noindent\begin{minipage}{0.55\linewidth}
    {\sffamily\footnotesize\color{lpGray}Lernpfad:}\\
    {\sffamily\small\color{lpInk}Wissen \textrightarrow{} Anwendung \textrightarrow{} Management}\\[10pt]
    {\sffamily\footnotesize\color{lpGray}Bearbeitungsdauer:}\\
    {\sffamily\small\color{lpInk}1 Kurstag}
  \end{minipage}\hfill
  \begin{minipage}{0.4\linewidth}
    \raggedleft
    {\sffamily\footnotesize\color{lpGray}Dozent:}\\
    {\sffamily\small\color{lpInk}Can Siebert}\\[10pt]
    {\sffamily\footnotesize\color{lpGray}Format:}\\
    {\sffamily\small\color{lpInk}Theorie \textperiodcentered{} Fallstudie \textperiodcentered{} Coaching}
  \end{minipage}

  \vspace{0.5cm}
  \noindent{\color{lpAccent}\rule{\linewidth}{0.6pt}}
\end{titlepage}

% --- LERNZIELE -----------------------------------------------------------
\section*{Lernziele dieses Tages}
\addcontentsline{toc}{section}{Lernziele dieses Tages}
\thispagestyle{plain}

\noindent Der letzte Kurstag schließt den Bogen: Regulatorik wie EU-Taxonomie und CSRD sind heute keine Reporting-Übung, sondern eine \emph{Prozess-, Daten- und Kontrollarchitektur}. Wer diese als Steuerungssystem versteht, integriert Compliance in die Transformation -- und macht aus dem Pflichtprogramm einen wirksamen Hebel für Datenqualität, Governance und Glaubwürdigkeit.

\subsection*{L1 -- Wissen (Reproduktion und Einordnung)}
\begin{itemize}
  \item EU-Taxonomie und CSRD als Datenbedarfe, Nachweise, Kontrollen und Verantwortlichkeiten verstehen -- nicht als Dokumentenpflicht.
  \item Governance- und Kontrollkonzepte einordnen: Policies, Controls, Audit-Readiness.
  \item Reportingprozess in 7 Schritten kennen: Plan -- Collect -- Validate -- Consolidate -- Approve -- Report -- Assure.
  \item Risiken erkennen: uneinheitliche Definitionen, Dateninseln, Schätzungen ohne Transparenz, fehlende Verantwortliche.
\end{itemize}

\subsection*{L2 -- Anwendung (Transfer auf konkrete Situationen)}
\begin{itemize}
  \item Einen End-to-End-Reportingprozess mit Rollen und Schnittstellen modellieren.
  \item Acht Controls entlang des Prozesses definieren (Vollständigkeit, Plausibilität, Freigabe, Version, Zugriff, Ausnahme).
  \item Ein Evidence Pack pro KPI strukturieren (Quelle, Berechnung, Freigabe, Änderungen, Ausnahmen).
  \item KPIs mit Anti-Gaming-Kopplung und Datenqualitätslabels designen.
\end{itemize}

\subsection*{L3 -- Management (Entscheidung und Verantwortung)}
\begin{itemize}
  \item Top-12-Entscheidungen einer ESG-Governance treffen (Systemgrenzen, Dictionary, Labels, Controls, Tooling, Cadence).
  \item Datenlücken im Dry Run transparent akzeptieren und mit Verbesserungsplan kompensieren.
  \item Transparenzentscheidungen (welche KPI bei Qualität C veröffentlicht?) bewusst treffen.
  \item Integration in Transformation sichern (Workflows, Process Intelligence, Green IT, Vendor).
\end{itemize}

\begin{merkbox}[Roter Faden]
EU-Taxonomie und CSRD sind keine ``Reporting-Projekte''. Sie sind \emph{Betriebssysteme} für Daten, Kontrollen und Entscheidungen -- mit dem zentralen Maßstab \emph{auditfähig}. Wer sie als Tabellen-Übung behandelt, baut Folien für die nächste Krise. Wer sie als Architektur baut, bekommt ein System, das gleichzeitig Compliance liefert und Steuerung ermöglicht. \quelle{vgl.\ EFRAG, 2023; European Commission, 2020}
\end{merkbox}

\clearpage

% --- 1. EINSTIEG ---------------------------------------------------------
\section{Einstieg: Vom Pflichtbericht zur Steuerungsarchitektur}

EU-Taxonomie und CSRD haben in den letzten Jahren die Regulatorik-Landschaft im Bereich Nachhaltigkeit grundlegend verändert. Was als Reporting-Pflicht für börsennotierte Unternehmen begann, weitet sich auf große mittelständische Unternehmen, deren Lieferketten und schließlich die gesamte Wirtschaft aus.

Die typische erste Reaktion: ``Wir müssen ein Reporting bauen.'' Die typische zweite, nach einigen Monaten: ``Wir müssen die Daten erst beschaffen, dann reporten.'' Die typische dritte, nach noch mehr Monaten: ``Wir müssen das in unsere Prozesse einbauen, sonst läuft das nie sauber.''

Die direkte Antwort spart Jahre: ESG-Reporting funktioniert nur, wenn es \emph{vom ersten Tag an als Architektur} verstanden wird.

\subsection{Drei typische Fehler}

\begin{enumerate}
  \item \fachbegriff{Dokumenten-Sicht statt Prozess-Sicht.} Excel-Dateien zwischen Werken, Einkauf und Finance herum-geschoben. Definitionen variieren, Zahlen widersprechen sich.
  \item \fachbegriff{Dateninseln ohne Transparenz.} Energieverbrauch im FM-System, Emissionsfaktoren im Excel, Lieferkette im Procurement-Tool. Niemand kann aggregieren.
  \item \fachbegriff{Schätzungen ohne Kennzeichnung.} Eine Zahl im Report wirkt präzise, ist aber aus Bauchwerten oder Branchen-Mitteln abgeleitet. Greenwashing-Risiko.
\end{enumerate}

\subsection{Was Regulatorik wirklich verlangt}

\begin{itemize}
  \item Klar definierte \fachbegriff{KPIs} mit \emph{Systemgrenzen}.
  \item \fachbegriff{Datenpipelines} von Quelle bis Report mit Validierung und Versionierung.
  \item \fachbegriff{Controls} (Plausibilität, Vollständigkeit, Freigabe).
  \item \fachbegriff{Evidence} (Nachweise zu Quelle, Berechnung, Freigabe).
  \item \fachbegriff{Governance} (Rollen, Entscheidungswege, Reviewzyklen).
\end{itemize}

\begin{eselsbox}[Eselsbrücke: \akzent{K-D-C-E-G}]
Fünf Bausteine einer ESG-Governance:\\[2pt]
\textbf{K}PIs (mit Systemgrenze)\\
\textbf{D}atenpipelines (Quelle bis Report)\\
\textbf{C}ontrols (Prüfungen entlang des Prozesses)\\
\textbf{E}vidence (Nachweise je KPI)\\
\textbf{G}overnance (Rollen, Entscheidungen, Reviews)\\[2pt]
\emph{K-D-C-E-G ist die Kurzformel für ``auditfähig''.}
\end{eselsbox}

% --- 2. EU-TAXONOMIE & CSRD ----------------------------------------------
\section{EU-Taxonomie und CSRD: das Mindset}

\subsection{EU-Taxonomie -- Klassifikation nachhaltiger Aktivitäten}

Die \fachbegriff{EU-Taxonomie} ist eine Klassifikation, die wirtschaftliche Aktivitäten danach beurteilt, ob sie zu sechs Umweltzielen wesentlich beitragen (Klimaschutz, Anpassung, Wasser, Kreislaufwirtschaft, Verschmutzung, Biodiversität) und gleichzeitig \emph{keine erhebliche Beeinträchtigung} anderer Ziele bewirken. \quelle{European Commission, 2020}

Für Unternehmen bedeutet das: bestimmte Kennzahlen (Anteil ``Taxonomie-fähiger'' Umsätze, CapEx, OpEx) müssen offengelegt werden. Das setzt eine systematische Erfassung der Aktivitäten und ihrer Klassifikation voraus.

\subsection{CSRD -- Corporate Sustainability Reporting Directive}

Die \fachbegriff{CSRD} verpflichtet Unternehmen ab einer bestimmten Größe (gestaffelt eingeführt) zur strukturierten Nachhaltigkeitsberichterstattung nach europäischen Standards (\fachbegriff{ESRS} -- European Sustainability Reporting Standards). Der Bericht ist Teil des Lageberichts und steht unter \fachbegriff{Limited Assurance} (perspektivisch Reasonable Assurance) -- d.\,h.\ wird durch externe Prüfer geprüft. \quelle{EFRAG, 2023}

Drei Konsequenzen:

\begin{itemize}
  \item Daten müssen genauso belastbar sein wie Finanzdaten -- Prozesse, Controls, Audit Trail.
  \item Definitions-Hoheit liegt nicht beim Unternehmen, sondern bei den Standards (ESRS).
  \item Transformationspläne müssen offengelegt werden -- nicht nur Zustand, sondern Pfad.
\end{itemize}

\subsection{Doppelte Wesentlichkeit}

Die CSRD verlangt eine \fachbegriff{doppelte Wesentlichkeitsanalyse} (\emph{Double Materiality}): \quelle{EFRAG, 2023}

\begin{itemize}
  \item \textbf{Impact Materiality.} Wie wirkt unser Unternehmen auf Umwelt und Gesellschaft?
  \item \textbf{Financial Materiality.} Wie wirken Nachhaltigkeitsfaktoren auf unser Unternehmen (Risiken, Chancen)?
\end{itemize}

Beide Perspektiven müssen analysiert und dokumentiert sein. Das ist nicht trivial -- aber es ist die Grundlage aller weiteren Reporting-Pflichten.

% --- 3. REPORTINGPROZESS -------------------------------------------------
\section{Reportingprozess: Plan -- Collect -- Validate -- Consolidate -- Approve -- Report -- Assure}

Ein belastbarer ESG-Reportingprozess hat sieben Schritte. Jeder Schritt hat eine Rolle und eine Kontrolle.

\subsection{Die sieben Schritte}

\begin{enumerate}
  \item \fachbegriff{Plan.} KPI-Dictionary aktualisieren, Datenanforderungen an Owner kommunizieren, Cadence festlegen.
  \item \fachbegriff{Collect.} Datenerfassung aus Quellen (IT, FM, Einkauf, HR, Finance). Mit Datenqualitätslabel.
  \item \fachbegriff{Validate.} Plausibilitätschecks, Vollständigkeitsprüfung, Sondersignale (Anomalien) markieren.
  \item \fachbegriff{Consolidate.} Aggregation auf Reporting-Ebenen (Werk, Bereich, Konzern).
  \item \fachbegriff{Approve.} Freigabe durch Owner und Reporting Lead. Zwei-Augen-Prinzip.
  \item \fachbegriff{Report.} Veröffentlichung im Lagebericht / Sustainability Statement / Dashboards.
  \item \fachbegriff{Assure.} Interne und externe Assurance (Limited / Reasonable). Sampling, Walkthroughs, Findings.
\end{enumerate}

\begin{figure}[h]
\centering
\begin{tikzpicture}[font=\sffamily\footnotesize,
  step/.style={rectangle, rounded corners=2pt, draw=lpAccent, fill=lpAccentLight,
               minimum width=1.55cm, minimum height=0.85cm, align=center, thick},
  last/.style={rectangle, rounded corners=2pt, draw=lpAccent, fill=lpGreenLight,
               minimum width=1.55cm, minimum height=0.85cm, align=center, thick},
  arr/.style={-{Stealth[length=2mm]}, lpGray, thick}]
  \node[step]                 (p1)               {Plan};
  \node[step, right=0.18cm of p1] (p2) {Collect};
  \node[step, right=0.18cm of p2] (p3) {Validate};
  \node[step, right=0.18cm of p3] (p4) {Consoli-\\date};
  \node[step, right=0.18cm of p4] (p5) {Approve};
  \node[step, right=0.18cm of p5] (p6) {Report};
  \node[last, right=0.18cm of p6] (p7) {Assure};
  \foreach \a/\b in {p1/p2,p2/p3,p3/p4,p4/p5,p5/p6,p6/p7} \draw[arr] (\a) -- (\b);
  \node[font=\sffamily\scriptsize\itshape, lpGray, below=0.15cm of p4]
    {jede Stufe: Rolle + Kontrolle};
\end{tikzpicture}
\caption{\sffamily ESG-Reporting-Pipeline -- sieben Schritte vom Plan bis zur externen Assurance.}
\end{figure}

\subsection{KPI-Dictionary als Fundament}

Ein \fachbegriff{KPI-Dictionary} ist die \emph{Single Source of Truth} für Definitionen. Pro KPI:

\begin{itemize}
  \item Name und ID.
  \item Definition.
  \item Systemgrenze (was ist drin, was draußen).
  \item Berechnungsformel.
  \item Datenquellen (Systeme, Owner).
  \item Datenqualitätslabel (A/B/C).
  \item Aktualisierungszyklus.
  \item Owner (Daten / Reporting / Approval).
\end{itemize}

Ohne KPI-Dictionary entstehen ``schwebende'' Definitionen -- jedes Werk versteht den KPI etwas anders, der Konzernbericht summiert Äpfel und Birnen.

\subsection{Annahmenregister}

Wo Schätzungen nötig sind (oft in Scope 3, Lieferkette), wird ein \fachbegriff{Annahmenregister} geführt:

\begin{itemize}
  \item Welche Annahme wurde getroffen?
  \item Auf welcher Grundlage?
  \item Wer hat sie freigegeben?
  \item Bis wann soll sie durch Messung ersetzt werden?
\end{itemize}

% --- 4. CONTROLS ---------------------------------------------------------
\section{Controls: prüfen, was zählt}

\fachbegriff{Controls} sind definierte Prüfungen entlang des Reportingprozesses. Sie sollten \emph{risikobasiert} sein -- nicht alles prüfen, sondern das Richtige. \quelle{ISACA, 2018}

\subsection{Acht typische Controls}

\begin{enumerate}
  \item \textbf{Completeness-Check.} Alle Pflichtfelder gefüllt? Alle relevanten Werke / Quellen liefern?
  \item \textbf{Plausibilitätsgrenzen.} Schwellenwerte für ``vermutlich falsch'' (z.\,B.\ Verbrauch $\pm$ 50\,\% gegen Vorjahr).
  \item \textbf{Zwei-Augen-Freigabe.} Pro KPI Owner + Reporting Lead bestätigen.
  \item \textbf{Änderungslog.} Wer hat was wann warum geändert?
  \item \textbf{Zugriffskontrolle.} Wer darf welche Daten sehen / ändern?
  \item \textbf{Ausnahmegenehmigung.} Wer darf von Standards abweichen, mit welcher Begründung?
  \item \textbf{Reconciliation zu Finance.} ESG-Kennzahlen mit Bezug zu Finanzdaten (z.\,B.\ Umsatz-Anteile) müssen abgleichbar sein.
  \item \textbf{Sampling-Review für Audit Readiness.} Stichprobenprüfung interner Art, simuliert das externe Audit.
\end{enumerate}

\subsection{Minimum Viable Controls vs.\ Overkill}

Im \emph{Dry Run} ist es legitim, eine reduzierte Control-Menge zu fahren -- mit dokumentierter Begründung und Verbesserungsplan. Welche Controls bleiben Pflicht:

\begin{itemize}
  \item Completeness (sonst fehlt die Basis).
  \item Plausibilität (sonst sind absurde Zahlen möglich).
  \item Freigabe (sonst ist unklar, wer verantwortet).
  \item Änderungslog (sonst ist nichts mehr nachvollziehbar).
\end{itemize}

\begin{merkbox}[Risk-based Controls]
Ein Control-Set, das alles prüft, prüft nichts mehr richtig. Risk-based bedeutet: Controls werden dort verdichtet, wo das Risiko hoch ist (Materialitäts-relevante KPIs, sensible Bereiche). Wo das Risiko gering ist, reicht ein leichteres Set. Diese Differenzierung ist Senior-Disziplin, nicht ``Kontrollvermeidung''.
\end{merkbox}

% --- 5. EVIDENCE PACK ----------------------------------------------------
\section{Evidence Pack: der Audit-Schlüssel}

Ein \fachbegriff{Evidence Pack} ist die Sammlung aller Nachweise, die ein Auditor zu einem KPI sehen wollen wird. Es ist nicht ``für den Audit'', sondern \emph{für die laufende Steuerung} -- der Audit ist nur die regelmäßige externe Prüfung dieser Sammlung.

\subsection{Sechs Pflicht-Artefakte je KPI}

\begin{enumerate}
  \item \textbf{Datenquelle.} Aus welchem System, welcher Export, welches Datum?
  \item \textbf{Transformationslogik.} ETL-Schritte, Skripts, Berechnungsformeln.
  \item \textbf{Berechnungsblatt / Script-Version.} Versionierte Berechnung, nachvollziehbar.
  \item \textbf{Freigabeprotokoll.} Wer hat wann auf welcher Basis freigegeben?
  \item \textbf{Ausnahmenregister.} Welche Abweichungen vom Standard, mit Begründung?
  \item \textbf{Datenqualitätslabel + Change-Requests.} Aktuelles Label, geplante Verbesserungsschritte.
\end{enumerate}

\subsection{Versionierung}

Jede Änderung (KPI-Definition, Datenquelle, Berechnung) erzeugt einen versionierten Eintrag. Auditoren wollen wissen: was war zum Reporting-Stichtag aktuell? Eine späte Korrektur ist erlaubt -- als nachvollziehbare Korrektur, nicht als ``stillschweigende Überschreibung''.

\subsection{Access Control für Evidence}

\begin{itemize}
  \item \textbf{Read-only.} Auditoren und Reviewer.
  \item \textbf{Append-only.} Daten-Lieferanten (sie schreiben Datensätze, ändern sie nicht).
  \item \textbf{Edit (mit Log).} KPI-Owner für Korrekturen mit Begründung.
  \item \textbf{Admin (selten).} Reporting Lead für strukturelle Änderungen.
\end{itemize}

% --- 6. KPI-DESIGN -------------------------------------------------------
\section{KPI-Design mit Anti-Gaming und Datenqualität}

Ein KPI ohne Anti-Gaming-Mechanik ist ein Anreiz zu Manipulation. Ein KPI ohne Datenqualitätslabel ist Greenwashing in Vorbereitung.

\subsection{Sechs bis acht ESG-KPIs als Set}

Beispielhaftes Set:

\begin{enumerate}
  \item \textbf{Energieverbrauch} (Klasse A/B/C nach Datenqualität).
  \item \textbf{Emissionsintensität} (mit Systemgrenze: Scope 1/2/3).
  \item \textbf{Anteil ``verifizierte Vendor-Daten''} (Lieferkette).
  \item \textbf{Datenqualitätsindex} (Aggregat über alle KPIs).
  \item \textbf{Control Pass Rate} (Anteil der Controls, die geprüft und bestanden wurden).
  \item \textbf{Timeliness} (Reporting pünktlich?).
  \item \textbf{Exception Rate} (Anteil KPIs mit Ausnahmen).
  \item Optional: \textbf{Transformationsfortschritt} (Anteil erreichter Ziele in Pfad).
\end{enumerate}

\subsection{Anti-Gaming-Kopplungen}

\begin{itemize}
  \item Kostenreduktion gilt nur, wenn Datenqualitätslabel $\geq$ B und Control Pass Rate $\geq$ X.
  \item Energie-Reduktion gilt nur, wenn Systemgrenze unverändert geblieben ist (sonst verschiebt sich der Erfolg in den ``draußen''-Bereich).
  \item Vendor-Verifikation gilt nur mit aktuellem Audit-Datum.
\end{itemize}

\subsection{Datenqualitäts-Labels}

\begin{itemize}
  \item \textbf{A -- gemessen.} Direkte Messung mit dokumentiertem Mess-Setup.
  \item \textbf{B -- berechnet.} Aus gemessenen Daten abgeleitet, mit dokumentierter Formel.
  \item \textbf{C -- geschätzt.} Auf Basis von Annahmen, Branchen-Mitteln, Proxy-Daten.
\end{itemize}

Pflicht: Jeder veröffentlichte KPI trägt ein Label. C-Werte erhalten zusätzlich einen Verbesserungsplan.

\begin{eselsbox}[Eselsbrücke: \akzent{A-B-C}]
Datenqualität auf einen Blick:\\[2pt]
\textbf{A} = gemessen (am verlässlichsten)\\
\textbf{B} = berechnet (aus A-Daten)\\
\textbf{C} = geschätzt (transparent + Verbesserungsplan)\\[2pt]
\emph{Eine Zahl ohne Label ist eine Behauptung. Eine Zahl mit Label ist Steuerungsinformation.}
\end{eselsbox}

% --- 7. FALLSTUDIE INDUSTRICO --------------------------------------------
\section{Fallstudie: IndustriCo Europe -- CSRD-Readiness}

\emph{IndustriCo Europe} betreibt mehrere Werke mit komplexer Lieferkette. ESG-Daten sind fragmentiert; einheitliche Definitionen fehlen. Reporting muss konsistent, nachvollziehbar und auditfähig werden. Gleichzeitig läuft ein Transformationsprogramm (Prozessdigitalisierung, Green IT, RPA).

Restriktionen: 14 Wochen bis zur ersten internen Dry-Run-Abgabe, Budget 250\,k EUR, Datenlücken (Energie, Lieferkette), Zielkonflikt Transparenz vs.\ Aufwand, Risiko ``Greenwashing''.

\subsection{Reportingprozess mit Rollen}

\begin{itemize}
  \item \textbf{Plan.} Reporting Lead + Process Owner aktualisieren KPI-Dictionary.
  \item \textbf{Collect.} Owner pro KPI in IT, Werken, Einkauf, HR, Finance.
  \item \textbf{Validate.} Datenteam mit Plausibilitäts-Tools.
  \item \textbf{Consolidate.} Reporting Factory (kleines Team).
  \item \textbf{Approve.} Owner + Reporting Lead (Zwei-Augen).
  \item \textbf{Report.} Veröffentlichung im Lagebericht und auf interner Plattform.
  \item \textbf{Assure.} Interne Sampling-Reviews, perspektivisch externe Limited Assurance.
\end{itemize}

\subsection{Acht Controls}

\begin{enumerate}
  \item Completeness-Check pro Werk.
  \item Plausibilitätsgrenzen (Energie $\pm$ 30\,\% gegen Vorjahr).
  \item Zwei-Augen-Freigabe je KPI.
  \item Änderungslog mit Begründungspflicht.
  \item Zugriffskontrolle role-based.
  \item Ausnahmegenehmigung mit Verbesserungsplan.
  \item Reconciliation gegen Finance-Daten (Umsatz, CapEx, OpEx).
  \item Sampling-Review (10\,\% der KPIs zufällig prüfen).
\end{enumerate}

\subsection{Evidence Pack je KPI}

\begin{itemize}
  \item Datenquelle (System / Export).
  \item Transformationslogik (ETL).
  \item Berechnungsblatt mit Version.
  \item Freigabeprotokoll.
  \item Ausnahmenregister.
  \item Datenqualitätslabel.
  \item Change-Requests.
\end{itemize}

\subsection{KPI-Set mit Anti-Gaming}

\begin{itemize}
  \item Energieverbrauch (A/B/C-Label).
  \item Emissionsintensität mit klarer Systemgrenze.
  \item Anteil verifizierte Vendor-Daten.
  \item Datenqualitätsindex.
  \item Control Pass Rate.
  \item Timeliness.
  \item Exception Rate.
\end{itemize}

Kopplung: Kostenreduktion / CO\textsubscript{2}-KPI nur gültig bei Label $\geq$ B und Control Pass Rate $\geq$ X.

\subsection{Integration in Transformation}

\begin{itemize}
  \item \textbf{DWH / ETL} für ESG-Events. Zentralisierung statt Excel-Inseln.
  \item \textbf{Workflow-Freigaben} für KPI-Änderungen. Versionierung und Audit Trail.
  \item \textbf{BI-Dashboards} mit eingebettetem Audit Trail.
  \item \textbf{RPA} für die Datensammlung (mit Governance aus Tag 12).
\end{itemize}

\subsection{Datenlücken im Dry Run}

\begin{itemize}
  \item \textbf{Lieferkettendaten} zunächst geschätzt (Label C) + Transparenz + Verbesserungsplan. Guardrail: Schätzung muss dokumentierte Methodik und Freigabe besitzen.
  \item \textbf{Energie in Außenstandorten}: Schätzung auf Basis Branchen-Mittel + Plan zur Messung in 6--12 Monaten.
\end{itemize}

% --- 8. COACHING ---------------------------------------------------------
\section{Coaching: Compliance durch Entscheidungen, nicht durch Folien}

Der häufigste Anti-Muster: Compliance als ``Erfüllen einer Vorgabe''. Die wirkungsvollere Haltung: Compliance als \emph{Entscheidungsarchitektur}.

\subsection{Drei Disziplinen}

\begin{enumerate}
  \item \textbf{Single Source of Truth.} KPI-Dictionary mit Systemgrenzen und Annahmenregister. Wenn jemand fragt ``Was meinst du mit X?'', gibt es eine eindeutige Antwort.
  \item \textbf{Risk-based Controls.} Materialität und Risiko bestimmen die Tiefe der Prüfung -- nicht ``alle gleich''.
  \item \textbf{Audit-Readiness im Alltag.} Versionierung, Freigaben, Evidence Packs sind keine Audit-Vorbereitung, sondern Alltagspraxis.
\end{enumerate}

\subsection{Senior-Reflexionsfragen}

\begin{itemize}
  \item Welche 3 Entscheidungen sind nicht delegierbar (Risk Appetite, Systemgrenzen, Transparenzregeln)?
  \item Wo entstehen die größten Reputationsrisiken -- und welche Controls verhindern sie?
  \item Welche KPI führt zu falschem Verhalten -- und wie wird sie gegen Gaming gekoppelt?
\end{itemize}

\subsection{Anti-Greenwashing-Test}

\begin{itemize}
  \item \textbf{Systemgrenzen klar benannt?}
  \item \textbf{Datenqualitätslabels sichtbar?}
  \item \textbf{Annahmen dokumentiert?}
  \item \textbf{Verbesserungsplan für C-Werte?}
  \item \textbf{Reconciliation zu Finance möglich?}
\end{itemize}

\begin{merkbox}[Anti-Greenwashing-Prinzip]
Die einfachste Anti-Greenwashing-Regel: \emph{Wenn ich diese Zahl nicht in einem Audit verteidigen könnte, veröffentliche ich sie nicht.} Wenn ich sie veröffentliche, aber sie ist nur geschätzt -- label sie. Wenn sie zu schwer messbar ist -- erkläre die Annahme. Transparenz ist die stärkste Verteidigung.
\end{merkbox}

% --- 9. MANAGEMENT-PERSPEKTIVE -------------------------------------------
\section{Management-Perspektive: ESG als Decision Architecture}

Auf Wissens-Ebene ist CSRD eine Verordnung. Auf Anwendungs-Ebene ist sie ein Reportingprozess. Auf Management-Ebene ist sie eine \emph{Decision Architecture} -- eine Sammlung von Entscheidungen, die jetzt getroffen werden müssen, weil sie sonst implizit getroffen werden.

\subsection{Top-12-Entscheidungen}

\begin{enumerate}
  \item Welche \textbf{Systemgrenzen} legen wir fest?
  \item Wer ist \textbf{Owner des KPI-Dictionary}?
  \item Welche \textbf{Datenqualitätslabels} verwenden wir, mit welcher Definition?
  \item Welches \textbf{Control-Set} ist Pflicht, welches optional?
  \item Welcher \textbf{Ausnahmeprozess} gilt?
  \item Welche \textbf{Tooling-Architektur} (DWH / ETL / BI / Audit Trail)?
  \item Welche \textbf{Reporting-Cadence} (monatlich, quartalsweise, jährlich)?
  \item Wie wird \textbf{Audit-Readiness} im Alltag gehalten?
  \item Welche \textbf{Vendor-Datenstrategie} verfolgen wir?
  \item Welche \textbf{Kommunikations- und Transparenzregeln} gelten?
  \item Welche \textbf{Priorisierung} (welcher KPI zuerst, welcher später)?
  \item Welches \textbf{Budget und welche Ressourcen} stehen bereit?
\end{enumerate}

\subsection{Governance und Operating Model}

\begin{itemize}
  \item \fachbegriff{ESG Steering.} Senior-Gremium, das strategische Entscheidungen trifft.
  \item \fachbegriff{Data Owners.} Pro Datendomäne (Energie, Personal, Lieferkette).
  \item \fachbegriff{Control Owners.} Pro Control-Familie (Plausibilität, Zugriff, Freigabe).
  \item \fachbegriff{Reporting Factory.} Operatives Team, das Konsolidierung und Reporting durchführt.
  \item \fachbegriff{Eskalation.} Zu wem geht ein Konflikt, eine Ausnahme, eine Krisensituation?
\end{itemize}

\subsection{Roadmap-Logik}

\begin{itemize}
  \item \textbf{Dry Run (4--6 Wochen).} Erste Daten-Erfassung, Minimum Viable Controls, Lessons Learned.
  \item \textbf{Stabilisierung (6--10 Wochen).} Lücken schließen, Controls erweitern, Tooling konsolidieren.
  \item \textbf{Assurance Readiness (laufend).} Audit-Sampling, externe Limited Assurance.
\end{itemize}

\subsection{Zielkonflikte}

\begin{enumerate}
  \item \textbf{Transparenz vs.\ Aufwand.} Mehr Detail kostet mehr Pflege -- aber zu wenig Detail erhöht Audit-Risiko.
  \item \textbf{Speed vs.\ Assurance.} Schnell veröffentlichen oder erst alle Controls grün?
  \item \textbf{Materialität vs.\ Vollständigkeit.} Auf das Wesentliche fokussieren oder alles abdecken?
\end{enumerate}

\begin{merkbox}[Schluss-Merksatz für Tag 16]
EU-Taxonomie und CSRD sind nicht das Ende der Transformation -- sie sind ihre \textbf{Steuerungsschicht}. Wer Compliance als Architektur baut, gewinnt nebenbei: Datenqualität, Klarheit, Governance, Glaubwürdigkeit. Wer Compliance als Pflichtübung erlebt, baut Folien für die nächste Krise. Senior-Verantwortung heißt: nicht ``erfüllen'', sondern \emph{steuern}.
\end{merkbox}

% --- 10. TAKE-AWAYS ------------------------------------------------------
\section{Take-aways aus Tag 16 (und dem Kurs)}

\begin{enumerate}
  \item \textbf{ESG ist Prozess, Daten, Controls -- nicht Dokument.} Wer Reporting als Tabellen-Übung baut, scheitert.
  \item \textbf{KPI-Dictionary als Single Source of Truth.} Mit Systemgrenze, Owner, Datenqualitätslabel.
  \item \textbf{Sieben-Schritt-Prozess}: Plan, Collect, Validate, Consolidate, Approve, Report, Assure.
  \item \textbf{Risk-based Controls}: Minimum Viable Controls im Dry Run, Erweiterung im Scale.
  \item \textbf{Evidence Pack je KPI}: Quelle, Berechnung, Freigabe, Änderungen, Ausnahmen, Datenqualitätslabel.
  \item \textbf{Anti-Gaming-Kopplungen}: Kennzahlen koppeln, damit Manipulation an einer Stelle die andere verschlechtert.
  \item \textbf{Doppelte Wesentlichkeit}: Impact und Financial Materiality beide analysieren.
  \item \textbf{Compliance als Decision Architecture}, nicht als Erfüllungspflicht. So entsteht Wirkung, nicht nur ein Bericht.
\end{enumerate}

% --- 11. LITERATUR -------------------------------------------------------
\clearpage
\section{Literatur}

Im Skript verwendete und für die Vertiefung empfohlene Quellen. Zitierstil: APA-7.

\begin{description}[style=nextline,leftmargin=18pt,labelindent=0pt,itemsep=4pt]
  \item[European Commission. (2020).] \emph{Regulation (EU) 2020/852 of the European Parliament and of the Council of 18 June 2020 on the establishment of a framework to facilitate sustainable investment (EU Taxonomy Regulation).} Official Journal of the European Union.
  \item[European Commission. (2022).] \emph{Directive (EU) 2022/2464 amending Directive 2013/34/EU as regards corporate sustainability reporting (CSRD).} Official Journal of the European Union.
  \item[European Financial Reporting Advisory Group (EFRAG). (2023).] \emph{European Sustainability Reporting Standards (ESRS) -- Set 1.} EFRAG.
  \item[International Organization for Standardization. (2015).] \emph{Environmental management systems -- Requirements with guidance for use} (ISO 14001:2015). ISO.
  \item[International Organization for Standardization. (2018).] \emph{Information technology -- Service management -- Part 1: Service management system requirements} (ISO/IEC 20000-1:2018). ISO.
  \item[ISACA. (2018).] \emph{COBIT 2019 framework: Governance and management objectives.} ISACA.
  \item[Project Management Institute. (2021).] \emph{A guide to the Project Management Body of Knowledge (PMBOK\textregistered{} Guide)} (7th ed.). Project Management Institute.
  \item[Westerman, G., Bonnet, D., \& McAfee, A. (2014).] \emph{Leading digital: Turning technology into business transformation.} Harvard Business Review Press.
\end{description}

% --- 12. GLOSSAR ---------------------------------------------------------
\clearpage
\section{Glossar}

Die wichtigsten Begriffe des Tages -- kurz definiert, in alphabetischer Reihenfolge.

\begin{description}[style=nextline,leftmargin=0pt,labelindent=0pt,itemsep=4pt]
  \item[Annahmenregister] Dokumentation aller Annahmen in KPI-Berechnungen (Grundlage, Freigeber, Verbesserungstermin).
  \item[Anti-Gaming (KPI)] Designprinzip: KPIs koppeln, sodass Manipulation an einer Stelle eine andere verschlechtert.
  \item[Assurance] Externe oder interne Prüfung der ESG-Berichterstattung. Limited Assurance (geringere Tiefe), Reasonable Assurance (höhere Tiefe).
  \item[Audit-Readiness] Alltagszustand der Auditfähigkeit: aktuelle Evidence Packs, Versionierung, Freigaben, Reconciliation.
  \item[Control] Definierte Prüfung im Reportingprozess (Vollständigkeit, Plausibilität, Freigabe).
  \item[CSRD] \emph{Corporate Sustainability Reporting Directive.} EU-Richtlinie zur Nachhaltigkeitsberichterstattung. \quelle{European Commission, 2022}
  \item[Datenqualitätslabel] A/B/C-Klassifikation der Datenqualität: A gemessen, B berechnet, C geschätzt.
  \item[Decision Architecture (ESG)] Sammlung der nicht-delegierbaren ESG-Entscheidungen einer Organisation.
  \item[Doppelte Wesentlichkeit] CSRD-Konzept: Impact (Wirkung auf Umwelt / Gesellschaft) und Financial Materiality (Wirkung auf Unternehmen). \quelle{EFRAG, 2023}
  \item[ESRS] \emph{European Sustainability Reporting Standards.} Konkretisierende Standards zur CSRD. \quelle{EFRAG, 2023}
  \item[EU-Taxonomie] Verordnung zur Klassifikation nachhaltiger wirtschaftlicher Aktivitäten gegen sechs Umweltziele. \quelle{European Commission, 2020}
  \item[Evidence Pack] Sammlung der Nachweise je KPI: Quelle, Berechnung, Freigabe, Änderungen, Ausnahmen.
  \item[KPI-Dictionary] Single Source of Truth für KPI-Definitionen mit Systemgrenze, Berechnung, Owner.
  \item[Materialität] Bewertung der Relevanz eines Themas für Reporting und Steuerung. CSRD verlangt doppelte Materialität.
  \item[Minimum Viable Controls] Kleinstes durchsetzbares Set an Controls; in Dry Runs zulässig, mit Verbesserungsplan.
  \item[Reconciliation] Abgleich von ESG-Kennzahlen mit Finanzdaten (z.\,B.\ Umsatz, CapEx, OpEx).
  \item[Reporting Factory] Operatives Team für Konsolidierung, Approval und Veröffentlichung der ESG-Daten.
  \item[Risk Appetite] Maß an Risiko, das die Organisation bewusst akzeptiert. Nicht delegierbar.
  \item[Scope 1 / 2 / 3] Klassifikation von Emissionen: direkt, eingekaufte Energie, Lieferkette und Nutzungsphase.
  \item[Systemgrenze] Festlegung, was in einer Bilanz oder Kennzahl berücksichtigt wird. Greenwashing-Risiko bei zu enger Wahl.
  \item[Timeliness] KPI-Eigenschaft: Reporting pünktlich abgeliefert?
  \item[Walkthrough] Audit-Technik: Schritt-für-Schritt-Durchgang eines Prozesses mit den Verantwortlichen.
\end{description}

\vspace{1cm}
\noindent{\color{lpAccent}\rule{\linewidth}{0.6pt}}\\[4pt]
{\footnotesize\sffamily\color{lpGray}Lernplatt \textperiodcentered{} by Can Siebert \textperiodcentered{} Kurs 71 (Dig 42) \textperiodcentered{} Tag 16 \textperiodcentered{} \textcopyright{} 2026}

\end{document}
